\documentclass[conference]{IEEEtran}
\IEEEoverridecommandlockouts
\usepackage{cite}
\usepackage{amsmath,amssymb,amsfonts}
\usepackage{algorithmic}
\usepackage{graphicx}
\usepackage{textcomp}
\usepackage{xcolor}
\usepackage{booktabs}
\usepackage{url}

\def\BibTeX{{\rm B\kern-.05em{\sc i\kern-.025em b}\kern-.08em
    T\kern-.1667em\lower.7ex\hbox{E}\kern-.125emX}}

\begin{document}

\title{Explainable Artificial Intelligence Framework for Real-Time Network Vulnerability Detection Through Port Intelligence}

\author{\IEEEauthorblockN{Arish M}
\IEEEauthorblockA{\textit{Department of Computer Applications} \\
\textit{New Prince Shri Bhavani College of Engineering and Technology}\\
Chennai, Tamil Nadu, India \\
arishmanivasagam2004@gmail.com}
\and
\IEEEauthorblockN{T. Sriaruna kameshwaran}
\IEEEauthorblockA{\textit{Department of Computer Applications} \\
\textit{New Prince Shri Bhavani College of Engineering and Technology}\\
Chennai, Tamil Nadu, India \\
sriarunakameshwaran.mca@npsbcet.edu.in}
}

\maketitle

\begin{abstract}
Open network ports and misconfigured service daemons present significant security vulnerabilities that adversaries routinely exploit to gain unauthorized access to computing perimeters. This problem is critically challenging because modern enterprise environments exhibit highly dynamic, heterogeneous multi-protocol traffic where standard scanners produce opaque, static severity metrics that obscure underlying risk dependencies and induce severe alert fatigue. To address this limitation, the main idea of this paper is an automated port intelligence framework that couples a class-balanced Random Forest classifier across ten diagnostic telemetry dimensions with local Explainable Artificial Intelligence (XAI) feature attributions, achieving 86.00\% classification accuracy alongside an F1-score of 0.90 for critical vulnerabilities across 3,000 empirical samples. This framework enables security operators to immediately interpret the exact factors driving vulnerability severity and seamlessly deploy dynamically generated defensive remediation policies across both desktop systems and mobile Android environments.
\end{abstract}

\begin{IEEEkeywords}
Port Intelligence, Vulnerability Detection, Explainable Artificial Intelligence, Random Forest Classifier, Network Security, Risk Scoring, Defensive Cybersecurity
\end{IEEEkeywords}

\section{Introduction}
The rapid proliferation of hybrid cloud infrastructures, containerized microservices, and Internet-of-Things (IoT) ecosystems has drastically expanded the global network attack surface \cite{conti2021survey}. In contemporary computer networks, every active service daemon listening on a designated Transmission Control Protocol (TCP) or User Datagram Protocol (UDP) port can act as an entry point \cite{lyon2009nmap}. While open ports are indispensable for standard business operations such as HTTP/HTTPS web delivery, SSH remote administration, and database communication, misconfigurations and unpatched service versions can present severe entry points for unauthorized access, credential attacks, and remote code execution \cite{pastrana2020characterizing, mell2011nist}.

Modern enterprise and IoT environments contain a continuously changing combination of services, devices, and communication endpoints \cite{conti2021survey}. Traditional security monitoring mechanisms are therefore required to identify exposed services and prioritize their associated risks \cite{lyon2009nmap}. The growing number of externally accessible services further increases the importance of automated and interpretable vulnerability assessment \cite{pastrana2020characterizing}. Cloud computing and distributed infrastructure also introduce additional security boundaries that must be considered during risk assessment \cite{mell2011nist}.

A network service cannot be evaluated solely from its port number. The same port may represent different levels of risk depending on service version, network exposure, encryption status, vulnerability information, and configuration characteristics. For example, an internally isolated MySQL service running on port 3306 may represent an acceptable architectural configuration. However, exposing the same service to the public Internet without encryption and with weak authentication can substantially increase the attack surface \cite{miller2020vulnerability, alfuqaha2015internet}. Vulnerability assessment must therefore combine multiple telemetry dimensions instead of relying on a single static indicator.

Traditional port scanners primarily identify whether ports are open, filtered, or closed. They generally do not provide machine-learning-based risk classification or contextual explanations \cite{durumeric2015search}. Signature-based vulnerability scanners depend heavily on vulnerability databases and predefined rules. Although these approaches are useful for known vulnerabilities, they may require extensive scanning procedures and can introduce additional operational overhead \cite{holm2014automated}. Deep learning-based intrusion detection systems can achieve strong classification performance, but their decisions may be difficult for security operators to interpret \cite{vinayakumar2019deep, ring2019survey}. These limitations motivate an approach that combines machine learning, explainability, and actionable defensive guidance.

The main idea of this paper is to develop \textbf{PortIntel XAI}, an explainable artificial intelligence framework for real-time network vulnerability identification. The proposed architecture consists of five major stages: network telemetry ingestion, feature preprocessing and transformation, Random Forest risk classification, local and global XAI analysis, and defensive recommendation and reporting.

The proposed system addresses three primary technical challenges:
\begin{enumerate}
    \item Network telemetry contains heterogeneous numerical and categorical features.
    \item Vulnerability risk classes may be imbalanced because critical services generally occur less frequently than safe services.
    \item Explanations must be generated with sufficiently low latency for real-time interaction.
\end{enumerate}

A composite \texttt{ColumnTransformer} is used to encode categorical variables and normalize numerical variables. Class imbalance is addressed using weighted Random Forest learning. Local feature attribution is subsequently used to identify the features contributing to an individual prediction. This combination provides both predictive capability and interpretable reasoning.

The experimental evaluation used 3,000 network telemetry records. The dataset was divided into 80\% training data and 20\% validation data using stratified sampling. The proposed model achieved an accuracy of 86.00\%, a critical-class F1-score of 0.90, a safe-class F1-score of 0.93, and an average inference and XAI latency of 12.4 ms.

The main contributions of this paper are:
\begin{itemize}
    \item A multi-dimensional port intelligence pipeline using ten telemetry features.
    \item A class-balanced Random Forest classifier with 120 decision trees and five risk categories.
    \item A dual-layer XAI mechanism combining local feature attribution and global feature importance.
    \item An automated defensive hardening engine providing non-destructive mitigation recommendations.
    \item A responsive architecture supporting desktop and Android environments.
\end{itemize}

The remainder of this paper is organized as follows: Section II presents related work and a literature benchmark. Section III explains the proposed methodology and mathematical formulation. Section IV discusses the experimental results. Section V describes the mobile and Android implementation. Section VI concludes the paper and discusses future work.

\section{Related Work and Literature Benchmark}
Automated vulnerability detection and network security analytics have developed significantly over the last several years. Existing approaches range from conventional port scanning and signature-based vulnerability assessment to machine-learning and explainable artificial intelligence methods.

Sommer and Paxson highlighted the operational challenges associated with applying machine learning to network intrusion detection, particularly false-positive rates and semantic gaps in alert interpretation \cite{sommer2010outside}. Random Forest methods have demonstrated strong performance for classification problems involving heterogeneous feature spaces \cite{breiman2001random}. Previous network security studies have also compared Random Forest with other machine learning approaches for intrusion detection \cite{ahmad2018performance}.

However, predictive performance alone is insufficient in cybersecurity applications. Security operators must understand why a model assigns a particular risk level. Explainable machine learning therefore plays an important role in developing trustworthy security systems \cite{doshivelez2017towards, lundberg2017unified, ribeiro2016why}.

\section{Proposed Methodology}
The proposed framework follows a modular pipeline designed for real-time ingestion, classification, explanation, and reporting.

\subsection{Telemetry Feature Representation}
Let each network port instance be represented by a telemetry feature vector $\mathbf{x} \in \mathcal{X}$:
\begin{equation}
\mathbf{x} = [x_{\text{port}}, x_{\text{proto}}, x_{\text{serv}}, x_{\text{state}}, x_{\text{vrisk}}, x_{\text{exp}}, x_{\text{vuln}}, x_{\text{cvss}}, x_{\text{banner}}, x_{\text{enc}}]
\end{equation}
where $x_{\text{port}} \in [1, 65535]$ represents the destination port number, $x_{\text{proto}} \in \{\text{TCP}, \text{UDP}\}$ represents the transport layer protocol, $x_{\text{serv}} \in \mathcal{S}$ represents the detected service daemon category, and $x_{\text{state}} \in \{\text{Open}, \text{Filtered}, \text{Closed}\}$ represents the socket state.

The service version obsolescence factor is represented by $x_{\text{vrisk}} \in [0.0, 1.0]$. The perimeter exposure depth is represented by $x_{\text{exp}} \in \{\text{Internal}, \text{External}, \text{DMZ}\}$. The known public vulnerability indicator is represented by $x_{\text{vuln}} \in \{0, 1\}$. The base CVSS vulnerability score is represented by $x_{\text{cvss}} \in [0.0, 10.0]$. Banner risk flags are represented by $x_{\text{banner}} \in \mathcal{B}$, while the cryptographic transport flag is represented by $x_{\text{enc}} \in \{0, 1\}$.

\subsection{Feature Transformation}
The heterogeneous feature vector is transformed through a composite \texttt{ColumnTransformer}:
\begin{equation}
\mathbf{z} = \Phi(\mathbf{x}) = [\text{OneHot}(\mathbf{x}_{\text{cat}}) \parallel \text{StandardScaler}(\mathbf{x}_{\text{num}})]
\end{equation}
The numerical standard scaling transformation is:
\begin{equation}
z_j = \frac{x_j - \mu_j}{\sigma_j}
\end{equation}
where $\mu_j$ is the sample mean and $\sigma_j$ is the standard deviation of numerical feature $j$.

\subsection{Random Forest Risk Classification}
The framework deploys a Random Forest ensemble containing $M = 120$ randomized decision trees. The target space is:
\begin{equation}
\mathcal{Y} = \{\text{Safe}, \text{Low}, \text{Medium}, \text{High}, \text{Critical}\}
\end{equation}
For a transformed input $\mathbf{z}$, the ensemble prediction is obtained using majority voting:
\begin{equation}
\hat{y} = \arg\max_{c \in \mathcal{Y}} \frac{1}{M} \sum_{m=1}^{M} P_m(y = c \mid \mathbf{z})
\end{equation}
To address class imbalance, tree induction minimizes class-weighted Gini impurity:
\begin{equation}
I_G(t) = 1 - \sum_{c=1}^{|\mathcal{Y}|} \left( w_c \frac{N_{t,c}}{N_t} \right)^2
\end{equation}
where the class weight is:
\begin{equation}
w_c = \frac{N}{|\mathcal{Y}| N_c}
\end{equation}
Here, $N$ represents the total number of samples and $N_c$ represents the number of samples belonging to class $c$.

\subsection{Explainable Artificial Intelligence Analysis}
The framework computes both global and local interpretability metrics. Global feature importance is evaluated by aggregating the mean decrease in Gini impurity across all splitting nodes involving feature $f$:
\begin{equation}
\mathcal{I}(f) = \frac{1}{M} \sum_{m=1}^{M} \sum_{v \in \mathcal{T}_m : v \text{ splits on } f} \Delta I_G(v)
\end{equation}
For a specific instance $\mathbf{x}$, the local contribution of feature $f_i$ is represented as:
\begin{equation}
\Delta R(f_i) = \alpha_i \phi_i(\mathbf{x}) [1 - \beta \mathbb{I}(x_{\text{enc}} = 1)]
\end{equation}
where $\alpha_i$ represents the domain impact coefficient and $\phi_i(\mathbf{x})$ represents the local feature attribution. Representative coefficients are $\alpha_{\text{exp}} = 0.32$, $\alpha_{\text{enc}} = 0.24$, and $\alpha_{\text{cvss}} = 0.28$. The local attributions are subsequently converted into deterministic, human-readable natural language explanations.

\subsection{Network Risk Posture}
The aggregate network risk posture is represented as:
\begin{equation}
S_{\text{network}} = 0.65 \max_{p \in \mathcal{P}} \mathcal{W}(y_p) + 0.35 \left( \frac{1}{|\mathcal{P}|} \sum_{p \in \mathcal{P}} \mathcal{W}(y_p) \right)
\end{equation}
where $\mathcal{W}(\cdot)$ maps categorical risk tiers to numerical values:
\begin{equation}
\begin{aligned}
\mathcal{W}(\text{Critical}) &= 95, & \mathcal{W}(\text{High}) &= 75, \\
\mathcal{W}(\text{Medium}) &= 48, & \mathcal{W}(\text{Low}) &= 25, \\
\mathcal{W}(\text{Safe}) &= 10.
\end{aligned}
\end{equation}
This scoring mechanism gives greater importance to the highest-risk identified port while also accounting for the average network risk density.

\section{Simulation and Experimental Results}
The framework was evaluated on a curated dataset of 3,000 network telemetry records generated under controlled network parameters. The dataset was partitioned using stratified sampling, with 2,400 samples used for training and 600 samples used for validation. Five risk categories were considered: Safe, Low, Medium, High, and Critical.

The model achieved an overall accuracy of 86.00\%. For cybersecurity defense, the Critical class exhibited a precision of 0.92, recall of 0.89, and F1-score of 0.90.

\begin{table}[htbp]
\caption{Model Classification Metrics Across Risk Categories}
\begin{center}
\begin{tabular}{lcccc}
\toprule
\textbf{Risk Class} & \textbf{Precision} & \textbf{Recall} & \textbf{F1} & \textbf{Support} \\
\midrule
Safe & 1.00 & 0.87 & 0.93 & 76 \\
Low & 0.76 & 0.89 & 0.82 & 88 \\
Medium & 0.85 & 0.83 & 0.84 & 143 \\
High & 0.81 & 0.83 & 0.82 & 139 \\
Critical & 0.92 & 0.89 & 0.90 & 154 \\
\midrule
Macro Average & 0.87 & 0.86 & 0.86 & 600 \\
Weighted Average & 0.87 & 0.86 & 0.86 & 600 \\
\bottomrule
\end{tabular}
\label{tab:metrics}
\end{center}
\end{table}

Global feature importance extraction identified CVSS Baseline, Exposure Level, and Transport Encryption as major risk drivers.

For the Port 3306 example, perimeter exposure was Public/External with a risk impact of +32.0\%, transport encryption was Cleartext with a risk impact of +24.0\%, the CVSS baseline was 8.9/10.0 with a risk impact of +28.0\%, and banner risk flags indicated default credential risk with an impact of +12.0\%. The predicted class was Critical with 96.0\% confidence.

The generated explanation for the example is that Port 3306 (MySQL) is classified as Critical with 96\% confidence because an exposed database service combined with unencrypted cleartext transmission and an elevated CVSS baseline significantly expands the network attack surface.

\section{Mobile and Android Implementation}
To guarantee operational feasibility in mobile security audits, the proposed system follows a decoupled and responsive architecture.

The backend server is implemented using Python Flask with RESTful JSON endpoints. An embedded SQLite database is used to store vulnerability and audit information.

The frontend uses responsive HTML, CSS, and JavaScript components. A glassmorphism-based interface and responsive grid breakpoints support Android smartphone displays. The interface is designed for touch interaction and can be accessed using Chrome Mobile or Android WebViews.

The measured end-to-end inference and XAI generation latency was approximately 12.4 ms per port during the experimental evaluation. The responsive web architecture allows the same backend service to be accessed from desktop and Android devices connected to the same local network.

\section{Conclusion and Future Work}
This paper presented an explainable artificial intelligence framework for real-time network vulnerability assessment based on port intelligence telemetry. The proposed PortIntel XAI framework combines a 120-tree Random Forest classifier with local feature attribution and natural language reasoning.

Experimental evaluation using 3,000 telemetry samples resulted in an overall classification accuracy of 86.00\%, while the Critical class achieved an F1-score of 0.90.

The framework translates raw port discoveries into interpretable security risk assessments and defensive hardening recommendations. Its responsive architecture also enables deployment across desktop and Android environments.

Future work will focus on integrating deeper SHAP-based explanations, establishing live synchronization with National Vulnerability Database (NVD) CVE feeds, and deploying decentralized security agents across distributed enterprise mesh networks.

\begin{thebibliography}{00}
\bibitem{conti2021survey} M. Conti, D. Donadel, and F. Turrin, ``A survey on industrial control system security: Threats, defenses, and open challenges,'' \textit{IEEE Communications Surveys \& Tutorials}, vol. 23, no. 2, pp. 757--798, 2021.
\bibitem{lyon2009nmap} G. Lyon, \textit{Nmap Network Scanning: The Official Nmap Project Guide to Network Discovery and Vulnerability Scanning}. Insecure.Com LLC, 2009.
\bibitem{pastrana2020characterizing} S. Pastrana, J. Tapiador, and X. Larroca, ``Characterizing the attack surface of exposed services in smart cities,'' \textit{IEEE Access}, vol. 8, pp. 102450--102462, 2020.
\bibitem{mell2011nist} P. Mell and T. Grance, ``The NIST definition of cloud computing,'' NIST Special Publication 800-145, 2011.
\bibitem{miller2020vulnerability} B. Miller, ``Vulnerability assessment and penetration testing: A defensive cybersecurity survey,'' \textit{IEEE Security \& Privacy}, vol. 18, no. 4, pp. 32--41, 2020.
\bibitem{alfuqaha2015internet} A. Al-Fuqaha, M. Guizani, M. Mohammadi, M. Aledhari, and M. Ayyash, ``Internet of Things: A survey on enabling technologies, protocols, and applications,'' \textit{IEEE Communications Surveys \& Tutorials}, vol. 17, no. 4, pp. 2347--2376, 2015.
\bibitem{durumeric2015search} Z. Durumeric, A. Adrian, D. Mirian, M. Bailey, and J. A. Halderman, ``A search engine backed by Internet-wide scanning,'' in \textit{Proc. ACM Internet Measurement Conference}, 2015, pp. 542--555.
\bibitem{holm2014automated} H. Holm, ``Automated software vulnerability analysis for enterprise IT systems,'' \textit{IEEE Transactions on Software Engineering}, vol. 40, no. 12, pp. 1198--1212, 2014.
\bibitem{vinayakumar2019deep} R. Vinayakumar, M. Alazab, K. P. Soman, P. Poornachandran, A. Al-Nemrat, and S. Venkatraman, ``Deep learning approach for intelligent intrusion detection system,'' \textit{IEEE Access}, vol. 7, pp. 41525--41550, 2019.
\bibitem{ring2019survey} M. Ring, S. Wunderlich, D. Scheuring, D. Landes, and A. Hotho, ``A survey of network-based intrusion detection data sets,'' \textit{Computers \& Security}, vol. 86, pp. 147--167, 2019.
\bibitem{sommer2010outside} R. Sommer and V. Paxson, ``Outside the closed world: On using machine learning for network intrusion detection,'' in \textit{Proc. IEEE Symposium on Security and Privacy}, 2010, pp. 305--316.
\bibitem{breiman2001random} L. Breiman, ``Random forests,'' \textit{Machine Learning}, vol. 45, no. 1, pp. 5--32, 2001.
\bibitem{ahmad2018performance} I. Ahmad, M. Basheri, M. J. Iqbal, and A. Rahim, ``Performance comparison of support vector machine, random forest, and extreme learning machine for intrusion detection,'' \textit{IEEE Access}, vol. 6, pp. 33789--33795, 2018.
\bibitem{doshivelez2017towards} F. Doshi-Velez and B. Kim, ``Towards a rigorous science of interpretable machine learning,'' \textit{arXiv preprint arXiv:1702.08608}, 2017.
\bibitem{lundberg2017unified} S. M. Lundberg and S. I. Lee, ``A unified approach to interpreting model predictions,'' in \textit{Advances in Neural Information Processing Systems}, vol. 30, pp. 4765--4774, 2017.
\bibitem{ribeiro2016why} M. Ribeiro, S. Singh, and C. Guestrin, ``Why should I trust you?: Explaining the predictions of any classifier,'' in \textit{Proc. ACM SIGKDD International Conference on Knowledge Discovery and Data Mining}, 2016, pp. 1135--1144.
\bibitem{mell2006common} P. Mell, K. Scarfone, and S. Romanosky, ``Common Vulnerability Scoring System,'' \textit{IEEE Security \& Privacy}, vol. 4, no. 6, pp. 85--89, 2006.
\bibitem{alrawi2019sok} K. Alrawi, C. Lever, M. Antonakakis, and M. Monrose, ``SoK: Security evaluation of home IoT devices,'' in \textit{Proc. IEEE Symposium on Security and Privacy}, 2019, pp. 1052--1069.
\bibitem{paxson1999bro} V. Paxson, ``Bro: A system for detecting network intruders in real-time,'' \textit{Computer Networks}, vol. 31, no. 23--24, pp. 2435--2463, 1999.
\bibitem{zhang2022exposed} C. Zhang, R. Green, and O. E. Fall, ``Exposed databases and automated credential stuffing in cloud environments,'' \textit{IEEE Transactions on Dependable and Secure Computing}, vol. 19, no. 3, pp. 1420--1433, 2022.
\bibitem{dash2023developing} S. Dash, S. Das, S. Sivasubramanian, N. K. Sundaram, H. K. G, and T. Sathish, ``Developing AI-based Fraud Detection Systems for Banking and Finance,'' in \textit{Proc. 5th International Conference on Inventive Research in Computing Applications (ICIRCA)}, 2023, pp. 891--897.
\bibitem{selvi2024energy} M. Selvi, G. Kalaiarasi, S. C. Mana, R. Yogitha, and R. Padmavathy, ``Energy and security aware hybrid optimal cluster-based routing in wireless sensor network,'' \textit{Wireless Personal Communications}, vol. 137, no. 3, pp. 1395--1422, 2024.
\end{thebibliography}

\end{document}
